% Independent arithmetic justification for the unchanged original 15671-witness circuit.
\subsection{The original forty-one-prime pattern}
Let
\[
 \mathcal D=\{0,1,2401,-343/32\},\quad R=47027=31\cdot37\cdot41,
 \quad W=25878772920=4\prod_{q\le29}q,
\]
\[
 c=18484837801,\qquad V_0=32W=828120733440.
\]
For integers $a_1,\ldots,a_k$, $k\le9$, put $A_i=Ra_i+3$,
$w_i=A_iv-u$. There are $5+4k$ displayed positions
\begin{equation}\label{eq:old41-list}
 v/V_0,\qquad u-\delta v,\qquad
 u-(A_i-\delta)v\quad(\delta\in\mathcal D,\ 1\le i\le k).
\end{equation}
For every lower bound these positions can simultaneously be positive
primes beyond that bound, congruent to $c$ modulo $W$, with the
following property. For the curve
$E:y^2=x^3-(3-\sqrt2)^2x$, both reductions above each prime are good,
and their cardinalities $d_1,d_2$ satisfy
\[
 \gcd(d_1,d_2)=4,\qquad p\nmid d_1d_2.
\]
This is a same-prime statement and asserts no cross-prime gcd condition.

\begin{proof}
Use $u=c+Wx$, $v=V_0(c+Wy)$. Put
\[
 \mathcal S=\mathcal D\cup\bigcup_i(A_i-\mathcal D),\quad
 K_s=V_0s\in W\mathbb Z,
 \quad\Psi_0=c+Wy,\quad\Psi_s=c+Wx-K_s(c+Wy).
\]
Remove duplicate targets. The remaining slopes are all distinct.
Indeed, the nonzero integral differences in $\mathcal D-\mathcal D$
are $\pm1,\pm2400,\pm2401$, none divisible by $31$.
The integral sums in $\mathcal D+\mathcal D$ are
$0,1,2,2401,2402,4802$, with residues $0,1,2,14,15,28$ modulo $31$,
none equal to $A_i\equiv3$. These facts exclude intersections between
secondary families and between primary and secondary families.
The $m=1+|\mathcal S|\le41$ forms have pairwise independent linear
parts. Repeated targets may be restored afterwards with repeated values.

At $q\mid W$, the forms equal $c$; direct reduction gives
$c\equiv1\pmod8$, $c\equiv3\pmod7$, and $c\equiv1\pmod q$ at
all other odd primes $q\le29$. Thus the forms and their successors
are units at every odd $q\mid W$. For $q>29$, the affine map from
$(x,y)$ to $(u,v)$ is invertible. At $q=31,37,41$, all $A_i=3$
modulo $q$, so $|\mathcal S\bmod q|\le8$. For $q\ge43$ this
cardinality is at most $40$. In either case choose distinct residues
$b,e$ outside the slopes. Put $\gamma=V_0^{-1}$ modulo $q$ and,
reversing $b,e$ if necessary, require $e-b\ne-\gamma$. Then
\[
 v=(e-b)^{-1},\quad u=bv
\]
satisfies
\[
 u-sv=v(b-s)\ne0,\quad u-sv+1=v(e-s)\ne0,
 \quad\gamma v\ne0,\quad\gamma v+1\ne0.
\]
This proves all the value and successor local conditions. No condition
on a successor is imposed at $2$.

Add an unused integer coordinate $z$, and choose a fixed rectangular
box $\Omega$ inside the cone $v>0$, $u/v>\max(0,|A_1|,\ldots,|A_k|)+2413$ in the
homogeneous leading coordinates. All forms on $X\Omega$ lie between
$\alpha X$ and $\beta X$ for positive fixed constants. Fix $Q>82$
containing every prime divisor of $W$ and every nonzero $K_s-K_t$.
Impose one good residue at every prime up to $Q$. For a cutoff $L>Q$,
exclude $\Psi_\nu=0,-1$ at every additional prime up to $L$.
The surviving proportion is
\[
 \gg_{\boldsymbol a}\prod_{Q<q\le L}(1-2m/q)
 \gg_{\boldsymbol a}(\log L)^{-2m},
\]
by \cite[Theorem~1, modulus~2]{williams-mertens}.

For fixed $L$, apply \cite[Main Theorem, p.~1763]{greentao} to every
surviving lattice residue in the convex box. The singleton partition
in \cite[Definition~1.5 and Lemma~1.6]{greentao} bounds complexity by
$m-2\le39$. The inputs $\mathrm{MN}(39)$ and $\mathrm{GI}(39)$ follow
from \cite[Theorem~1.1]{moebius} and \cite[Theorem~1.3]{gtz}.
At imposed primes the local factor is $(1-1/q)^{-m}$; at the other
primes it is $1+O(m^2/q^2)$. The singular product therefore has a
positive lower bound independent of $L$. All coefficient and modulus
bounds are fixed before $X$ tends to infinity. Summing the finitely
many asymptotics and removing proper prime powers gives
\begin{equation}\label{eq:old41-lower}
 \#\{\text{prime patterns}\}\ge
 c_{\boldsymbol a}\frac{X^3}{(\log L)^{2m}(\log X)^m}.
\end{equation}
The constant $c_{\boldsymbol a}>0$ is independent of $L$; the threshold
for $X$ may depend on it. Proper powers contribute at most
$O_{\boldsymbol a}(X^{5/2}(\log X)^{m+1})$ before removing von Mangoldt
weights, since each fixed value has $O_{\boldsymbol a}(X^2)$ fibers.

We give the uniform upper estimates needed to impose the order condition.
For forbidden sets of cardinality $\rho(q)$ in a two-dimensional box
of side $N$, put $g(q)=\rho(q)/q^2$, $h(q)=g(q)/(1-g(q))$, and extend
multiplicatively on squarefree support. CRT gives
\[
 A_d=N^2g(d)+R_d,\qquad |R_d|\ll Nd+d^2.
\]
With $G(Z)=\sum_{d\le Z}h(d)$, take
\[
 \lambda_d=\frac{\mu(d)}{G(Z)\prod_{q\mid d}(1-g(q))}
 \sum_{\substack{e\le Z/d\\(e,d)=1}}h(e)\quad(d\le Z).
\]
Then $\lambda_1=1$, $|\lambda_d|\le d$, and the Selberg quadratic
form equals $G(Z)^{-1}$. To verify the last assertion, use
$y_r=\sum_{r\mid d}\lambda_dg(d)=\mu(r)h(r)/G(Z)$ and
$1/g(n)=\sum_{r\mid n}1/h(r)$. The number of sifted points is thus
\begin{equation}\label{eq:old41-upper}
 \ll N^2/G(Z)+NZ^6+Z^8.
\end{equation}
For $\rho(q)=2q-1$, supported on $q\equiv1\pmod4$ with at most one
prime omitted, $G(Z)\gg_Q\log Z$, uniformly in that omission.
For $\rho(q)=(m-1)q$ at all $q>Q$,
$G(Z)\gg_{Q,m}(\log Z)^{m-1}$. These follow from
\cite[Theorem~1, moduli~4 and~2]{williams-mertens}. Truncation is
justified by using $z=Z^{1/512}$ and the probability distribution
$h(d)/\prod_{q\le z}(1+h(q))$: its expected $\log d$ is at most
$40\sum_{q\le z}(\log q)/q<\tfrac12\log Z$ for large $Z$.
Thus at least half the mass lies below $Z$. Removing one norm-sieve
prime changes its product by a factor at least $1/2$.

For a selected prime, $(7/p)=-1$ and $p\equiv1\pmod8$. The two
reductions of $3-\sqrt2$ have product $7$ and opposite square classes,
so their traces are opposite. Their discriminants $64d^6$ are units.
The automorphism $(x,y)\mapsto(-x,iy)$ is rational over $\mathbb F_p$.
By the Tate-module injection \cite[III.7.4]{silverman}, its rational
centralizer is $\mathbb Q(i)$: at the auxiliary prime $3$ the
centralizer has dimension two because $T^2+1$ is irreducible.
Frobenius is integral by \cite[V.2.3.1(b)]{silverman}; hence it is
$a+bi\in\mathbb Z[i]$. The trace and determinant formulas
\cite[III.8.6]{silverman} give
\[
 a^2+b^2=p,\qquad d_1=p+1-2a,\qquad d_2=p+1+2a.
\]
Full rational two-torsion gives $4\mid d_1,d_2$, and
$v_2(d_1+d_2)=2$, so their gcd has $2$-part exactly $4$.
Also $0<d_i<2p$ and $d_i\ne p$, so $p\nmid d_1d_2$.
An odd common divisor $\ell$ divides $a$ and $p+1$, and forces
$b^2\equiv-1\pmod\ell$. The successor conditions exclude $\ell\le L$.

The number of bad prime values of any fixed form is
\begin{equation}\label{eq:old41-bad}
 \ll_{\boldsymbol a}X/(L\log X)+X^{7/8}+\sqrt X\log X.
\end{equation}
Indeed, for $L<\ell\le X^{1/8}$, write
$a=\ell r$, $b=b_0+\ell s$, with at most two $b_0$ modulo $\ell$.
The box has side $N\asymp\sqrt X/\ell\gg X^{3/8}$.
Use the norm case of \eqref{eq:old41-upper} with $Z=X^{1/100}$;
its error is uniformly $o(N^2/\log X)$. Sum
$O(X/(\ell^2\log X))$. For $\ell>X^{1/8}$, $a\ne0$ gives
$O(\sqrt X/\ell)$ choices with no additive constant, and the
$b$ count is $O(\sqrt X/\ell+1)$. Summation proves the remaining terms.

On a fiber $\Psi_0=p$, use $(x,z)$; the other forms are
$c+Wx-K_s p$. On a fiber $\Psi_{s_0}=p$, use $(y,z)$ and
\[
 x=K_{s_0}y+\frac{p-c(1-K_{s_0})}{W},\qquad
 \Psi_s=p+(K_{s_0}-K_s)\Psi_0.
\]
The quotient is integral for an actual fiber. For
$Q<q\le X^{1/100}<p$, the remaining zero lines are distinct and
parallel, with total cardinality $(m-1)q$. The second case of
\eqref{eq:old41-upper} bounds each fiber by
$O_{\boldsymbol a}(X^2/(\log X)^{m-1})$, independently of $p,L$.
Multiply by \eqref{eq:old41-bad} and sum over $m$ positions.
The bad patterns are at most
\[
 C_{\boldsymbol a}\frac{X^3}{L(\log X)^m}
 +O_{\boldsymbol a}\left(
 X^{23/8}/(\log X)^{m-1}+X^{5/2}/(\log X)^{m-2}\right).
\]
Choose $L$ so large that $C_{\boldsymbol a}/L$ is less than half
$c_{\boldsymbol a}/(\log L)^{2m}$, then choose $X$ large.
A good pattern survives, proving the assertion.
\end{proof}

\subsection{A shared auxiliary point with private shifts}
Let $\mathcal P$ be the finite set of distinct selected primes, all
larger than $76176$. At each $p$ fix an orientation of its two places,
and let $d_{p,1},d_{p,2}$ be the corresponding group cardinalities.
Set
\[
 j=\operatorname{lcm}_{p\in\mathcal P}d_{p,2}>0,
 \qquad J=[j]P_*.
\]
For each $p$, independently choose an integer $h_p$ satisfying
\begin{equation}\label{eq:old41-private-crt}
 h_p\equiv0\pmod{d_{p,1}},\qquad
 h_p\equiv4\pmod{d_{p,2}},\qquad
 h_p\equiv4\pmod{p^{64+v_p(j)}}.
\end{equation}
The first two congruences are compatible because their modulus gcd
is $4$. The last is compatible because $p\nmid d_{p,1}d_{p,2}$.
The resulting arithmetic progression contains an $h_p$ outside
$\{0,4,-4\}$. Put
\[
 X_0=[h_p]P_*,\qquad X_i=[h_p n_i]P_*,\qquad
 Y_i=[(h_p-4)n_i]P_*,\quad n_0=1.
\]
These are private for each contact occurrence, whereas $J$ is common.
Use the same orientation at repeated occurrences of $p$.

At the primary place all $X_i$ are formal; at the opposite place
$J$ and all $Y_i$ are formal. The formal logarithm is a homomorphism
and an isometry on $p\mathbb Z_p$ by
\cite[IV.6.4 and its proof; VII.2.1--2.2]{silverman}. The hypotheses hold
because $p$ is odd, split, and of good reduction. For any integers
$k,j$ whose multiples of $P_*$ lie in the formal group,
\[
 \frac{\log([k]P_*)}{\log([j]P_*)}=\frac{k}{j}.
\]
To see this without defining a logarithm of $P_*$ itself, apply the
homomorphism law to $[j]([k]P_*)=[k]([j]P_*)$. Nonzero formal
multiples have nonzero logarithm. Therefore
\[
 v_p\left(\frac{z(Y_i)}{z(J)}\right)
 =v_p\left(\frac{(h_p-4)n_i}{j}\right)\ge64,
 \qquad
 \frac{\log(X_i)}{\log(X_0)}=n_i.
\]
This proves the required local comparison conditions. Although other
$d_{q,2}$ may be divisible by $p$, their effect is exactly the term
$v_p(j)$ in \eqref{eq:old41-private-crt}; no cross-prime gcd assertion
is needed.

For the original truncation
\[
 L(z)=z\sum_{r=0}^{15}\frac{(-1)^r\binom{2r}{r}}{4r+1}(Dz^4)^r,
\]
the full logarithm has the same series extended over all $r\ge0$.
This follows by integrating
$(1+4Dz^4)^{-1/2}$, the invariant differential in the parameter
$z=-x/y$. If $v_p(z)\ge1$ and $p>76176$, every omitted term has
valuation at least $65v_p(z)$: its integer numerator can only increase
valuation, and for $n=4r+1\ge65$,
$n v_p(z)-v_p(n)\ge65v_p(z)$.
For $n<p$ this is immediate; for $n\ge p$ use $n-65\ge v_p(n)$.
Thus $L(z)/\log(z)\in1+p^{64}\mathbb Z_p$. The required primary
truncated-log ratios are congruent to the integers $n_i$ modulo $p^{64}$.

Conversely, in an accepted comparison at any $p>76176$, its conjugate
parameter-ratio tests imply
\[
 m_i-4n_i\equiv0\pmod{p^{64}},\qquad h-4\equiv0\pmod{p^{64}},
\]
since $j$ is a nonzero integer. Hence $h$ is a $p$-unit. The same
logarithm identity and truncation estimate give
$L(X_i)/L(X_0)\equiv m_i/h\equiv n_i\pmod{p^{64}}$ and show
that this ratio is integral. The rational proxy test then gives the
claimed integrality and congruence for every $b_i$.

\subsection{The common root approximation and parity boundary}
For each distinct selected $p$, retain the chosen square root
$\rho_p$ of $2$. After constructing the private points, each field
test $A+B\sqrt2$ of order $k$ has $v_p(A+B\rho_p)\ge k$.
Choose a common integer $\theta$ sufficiently close to $\rho_p$ at
every selected prime that
$v_p(\theta-\rho_p)\ge k-v_p(B)$ for every nonzero coefficient $B$,
and at least one. There are finitely many requirements; CRT supplies
such an integer. Since $\theta+\rho_p$ is a unit,
\[
 v_p((\theta^2-2)B)\ge k,\qquad v_p(A+\theta B)\ge k.
\]
Repeated primes require the same root and impose compatible precision
bounds. Conversely, the root test and these two rational tests imply
the field test at the selected root. This verifies the circuit's one
common $\theta$, not a separate root variable per contact.

Every raw contact parameter in the complete assignment is a positive
prime congruent to $c$ modulo $W$. It is therefore a square in
$\mathbb Q_2,\mathbb Q_3,\mathbb Q_{17}$. The binary norm batch has
unit root at its contact prime, because normalized leaves are integral
and its leftmost constant leaf is $1$. Its root is positive and nonzero
in the real embedding. The argument $r\mathcal W^2$ consequently is
positive, has square reduced denominator, and is locally square at
$2,3,17$. These are the exact hypotheses of the original specialized
parity-completeness input. This verifies its consumption but does not
constitute a new proof of that input. In particular no integrality of
the comparison points at these three small primes is required here.
